%% ===== Abstract =====
% [27/10] Viết lại cả abstract sau review: số theo phương án B (Compilation 71%, O-e và O-h 4%);
% bỏ vế "only in 2026" (V5); "relies on human judgment" thành "9 of the 12 ... include a human check"
% (comment 1: chỉ 3/9 bài LIB có người chấm mà không chạy code); thêm pilot, 52% đến 90% (P1);
% tách câu dài để mỗi dấu câu chỉ một chức năng.

% [29/10] "Nine" thành "Ten" theo Bảng 10 (Google migrations mã lại); "six gaps" thành "five gaps" theo Discussion (G1--G5).
% [30/10] Thêm nguồn test theo ngữ cảnh (RQ2, cơ sở G3): "existing tests" làm chủ ngữ cả hai vế vì chỉ 7/14 bài LIB, VER dùng test sẵn có (6/7 lấy từ suite của dự án); XL 12/20 benchmark. "such tests" thay "them" cho khỏi lẫn.
% [30/10] Thêm 14/49 = 29% chạy bản gốc (G5); "Nearly a fifth" thành 18% (9/49). "report separate results" vì 40 bài kết hợp oracle, 33 bài báo kết quả riêng (như Results). Chênh compile/behavior như Conclusion (13/18 lệch hơn 10 điểm); "Yet" nối câu 10 bài không nêu hạn chế.
% [30/10] "papers on legacy migration" như Conclusion; "Across oracles and thresholds" vì biên dịch và test không có ngưỡng; bỏ "even"; câu pilot thì quá khứ ("varied", "judged"); "five gaps…" theo RQ5; "our data and code" như Intro.
\begin{abstract}\cmtpink{P01, Major; EIC, DA; open since R01}{Pink notes are the 30/10 review by a five-reviewer panel plus deep research; reasons, evidence and the revision plan are in Review\_Claude\_Seminar\_2\_3010.md at the repository root. Panel decision: major revision. Before any clean build: setting reviewfalse in macros.tex still fails (Undefined control sequence hcrev, about 100 errors), because hcrev and revtag are defined only in the review branch; add pass-through definitions to the else branch. Both builds still print the placeholder author block.}
\new{Large language models (LLMs) are now used to translate code between languages and to migrate it across versions, libraries and frameworks.\cmtpink{P02, Critical for EASE; EIC}{The EASE build is 18 pages (NeurIPS build 28) against a 10-page limit; over-length papers are desk-rejected. Cut plan (EIC-01 in the report): move the Appendix table and Tables 2, 4 and 6 and the track and year rows of Table 8 to the replication package; Discussion from about 3.4k to 1.3k words; pilot to about 800 words; this abstract from about 340 words and 20 numbers to 200--250 words. If it still does not fit, publish the pilot separately.} Whether the generated code behaves like the original is decided by an oracle, and the choice of oracle can change the reported result. We present a systematic mapping study of how behavioral equivalence is evaluated in LLM-based code translation and migration,\cmtpink{P03, Minor; EIC, R2}{The abstract never says what is new compared with the existing reviews (Chen et al., Zhang et al., Rgaguena et al.). Add one clause, e.g. prior reviews catalogue translation metrics but do not trace where test ground truth comes from or how far oracles disagree on the same outputs.} covering the main proceedings of ICSE, FSE and ASE from 2024 to 2026.\cmtpink{P04, Minor; R1, DR; open since R02, R10}{Main proceedings here also covers Companion-volume papers (some from workshops), posters and journal-first talks (Section 4.1), and refs.bib gives seven FSE Companion papers main-proceedings booktitles. Most readers will take it to mean the research-track volume; say proceedings (main and Companion volumes).} From 4,451 papers we selected 49 and coded each from its full text. The most common oracles are compilation (71\%) and existing tests (57\%); formal verification, LLM-generated tests and LLM judges are rare (4\% each).\cmtpink{P05, Major; R2, DA, R1; open since R08}{Two qualifiers. (1) Compilation 71\% counts any reported compile count, including failure breakdowns (Section 4.5), and 33 of these 35 papers also run the code (all but CodePlan and DepUpgrade), so compilation is almost never the only evidence; say so. (2) LLM-generated tests 4\% undercounts LLM-written tests: O-e is defined by where expected outputs come from, and 3 more papers run LLM-generated inputs under O-f. See the Table 1 comment.} Only 29\% of the papers run the original program to compare behavior\cmtpink{P06, Major; DA, R2}{A project suite that passed on the pre-migration code already compares against the original's behavior on its inputs, so 29\% run the original is not the same as 71\% never compare with the original. Reword (e.g. only 29\% run the original program on the same inputs as the migrated code) and see the G5 comment.}, and 18\% use no oracle that executes the code. In translation between languages, existing tests come mostly from benchmarks; in library and version migration, they usually come from the project's own suite. None of the 5 papers on legacy migration uses such tests\cmtpink{P07, Major; EIC, R2, R1, DA}{This depends on a legacy flag that excludes C, although at least 4 mapped C-to-Rust papers call their sources legacy and use existing tests; all 5 papers are also industry or NIER papers, and legacy systems often have no automated tests. Drop it from the abstract or qualify it (5 papers with COBOL, PL/I or mainframe sources). See the comment on the legacy flag in Section 2.}, and 9 of the 12 on library migration include a human check. Thirty-three papers report separate results from two or more oracles on the same outputs. In 13 of the 18 that give both a compile rate and a behavioral rate, the compile rate is more than 10 percentage points higher.\cmtpink{P08, Major; R1, DA, R2, EIC; open since R21}{13 of 18 depends on which configuration represents each paper (no rule is stated; 4 rows are baselines, an ablation or pre-repair output) and on a 10-point cut that is not justified. See the comment in Section 5.4; also give the median gap (19 points).} Yet 10 of the 49 papers state no limitation of their own oracle.\cmtpink{P09, Minor; DA, EIC}{5 of these 10 are short papers (IVR, two posters, a demo, NIER), i.e. 5 of the 7 short papers against 5 of the 42 full papers, so page limits plausibly explain half. Yet also links two unrelated findings. Drop Yet and add half of them short papers; note it in Threats and next to checklist item 7.} In a pilot that varied only the oracle, we judged an LLM's translations of 164 HumanEval-X tasks from Python to Java with four oracles: compilation, the benchmark tests, CodeBLEU and a second LLM as judge. Across oracles and thresholds, 52\% to 90\% of the same translations count as correct.\cmtpink{P10, Major; EIC, R1, DA; open since R04}{Both ends are weak anchors: 90\% is compilation, which no mapped paper treats as correctness, and 52\% is CodeBLEU at 0.7, a cut-off you chose and do not treat as main (67\% at 0.6). The mapped papers report CodeBLEU as an average, not as a pass/fail verdict. Suggest: between 67\% (CodeBLEU at the main threshold) and 90\% (compilation) of the same translations are accepted, against 80\% under the tests; and name the setup (a 7B translator, one run). The same sentence is in the Introduction and the Conclusion.} The judge accepts half of those that fail the tests and rejects 41\% of the benchmark's own reference solutions.\cmtpink{P11, Major; DA, R1, R2, EIC; open since R04}{Say a small (8B, 4-bit) LLM judge, as the Conclusion does. The 41\% is not yet judge error: the references are not proven equivalent to the Python code and none of the 67 rejected ones was inspected; a judge asked about every valid input can rightly reject, e.g., an int-overflow difference. See the pilot comments.} A success rate is thus meaningful only together with the oracle that produced it. From these results we derive five gaps between current evaluation practice and the needs of real migration projects\cmtpink{P12, Major; DA, R3, EIC, R1}{The needs side rests on one premise argued in Section 6 (a migration replaces an existing program, so it needs evidence...). No practitioner data, acceptance criteria or V\&V standard is analysed, and the map's own industry papers partly cut against it. Reword here, in RQ5, contribution (5) and Section 6, e.g. gaps between how published studies evaluate migrations and the evidence that behavioral equivalence requires; the Conclusion's what a migration needs is already closer.}, along with a checklist for reporting oracle provenance. We release our data and code.\cmtpink{P13, Major; EIC, R1}{The release tag is still a placeholder and, per the README, the repository has no release or license yet. Several printed counts come from reading rather than from coded fields (18 papers classify failures, 5 describe faults that slip past one check, the G4 cases), so the package cannot recompute them until reading\_counts.csv is added. Create a tagged, licensed release archived with a DOI before submission; for EASE use an anonymized mirror.}}
\end{abstract}